\chapter{コンプライアンス・説明責任}

\section{背景：専門家の判断を社会の信頼へつなぐ}
技術者は安全、仕様、費用について依頼者より多くの情報を持つ場合がある。製品の使用者や施設周辺の住民は、その判断の結果を受ける一方で、判断の過程を直接見ることが難しい。適切な計算だけでなく、根拠を記録し、他者の権利を尊重し、判断を説明する仕組みが必要になる。

法令遵守は法律に触れない範囲を知るだけではない。誰のどの利益を保護する制度かを理解し、現実の行為を要件に照らすことである。説明責任、競争の公正、知的創作への報酬、安全上の問題の通報、規格への適合は、異なる仕組みを通じて信頼を支える。本章は原典確認した問題の背景初稿であり、未処理の全問題への執筆を継続する。

\section{説明責任：情報を出すことと理解できること}
\label{theory:compliance:accountability}
説明責任とは、自分の判断の根拠、結果、制約について、関係者へ説明し吟味に応じる責任である。報告書を公開した事実だけで果たせるとは限らない。何を測り、何を仮定し、どこに不確実さがあり、どの代替案と比較したかを、相手が検討できる形で伝える。

専門外の公衆に専門家と同じ知識を要求するのでなく、用語・図・尺度を調整して理解を支援する。説明への質問や異なる価値判断を、知識不足という理由だけで退けない。技術士倫理綱領は真実性と秘密の適切な扱いを示し、電気学会行動規範も相手の専門性に応じた誠実な議論を求める\parencite{ipej_ethics2011,ieej_conduct2007}。

説明責任と秘密保持はともに必要である。第三者の個人情報や無関係な技術秘密を全公開するのでなく、必要な情報・範囲・相手を検討する。安全上の重要情報を秘密として隠すことも許されない。危険の程度、緊急性、通報先、開示の必要性を記録して判断する。

\section{合意形成：事実の共有と公正な手続}
\label{theory:compliance:consensus}
合意形成では、事業者、利用者、住民などの利益と価値観が異なる。騒音を何dBと測るかは技術的問いであるが、誰がどの負担を引き受けるかは測定だけで決まらない。事実の確認と利害の調整を分けることが出発点である。

関係者分析により、参加すべき人、懸念、守りたい利益を把握する。進行役は発言機会と議論の手順を整え、技術助言者は専門的な検討を支える。担当技術者は専門能力を持っていても、事業の当事者として利害を持ち得る。中立性と進行能力は別に確保する必要がある。国土交通省の2006年研究と住民参加資料は、第三者の中立性、関係者分析、専門助言と手続支援を扱う\parencite{mlit_mediation2006,mlit_pi_thirdparty}。第三者の参加だけで必ず合意が成立するとは限らず、不同意の理由も記録する。

\section{知的財産：技術、表現、標識、秘密の保護}
\label{theory:compliance:ip}
研究開発や創作には費用がかかるが、成果を得た後の模倣は比較的容易なことがある。制度は創作への誘因と成果の利用を調整する。保護対象が異なるため、全てを特許と呼ぶことはできない。

\begin{center}
\begin{tabular}{p{.19\linewidth}p{.69\linewidth}}
\hline
権利・利益 & 保護対象の基本的な違い\\ \hline
特許 & 自然法則を利用した技術的な発明。新規性・進歩性等は別途審査する。\\
実用新案 & 物品の形状・構造・組合せに関する考案。方法一般の保護ではない。\\
意匠 & 物品等の外観に関するデザイン。具体的な保護対象は法の時点に従う。\\
商標 & 商品・サービスの出所を識別する標識。技術的仕組みとは異なる。\\
著作権 & 思想・感情の創作的な表現。アイデアだけの独占ではない。\\
営業秘密 & 要件を満たして秘密管理する有用な非公知の技術・営業情報。\\ \hline
\end{tabular}
\end{center}
知的財産基本法の範囲は広く、特許・実用新案・意匠・商標の4つが産業財産権である。著作権も知的財産権であるが、この4種類には含まれない\parencite{ip_basic_act,jpo_ip_overview}。平成28年・令和4年・令和5年の設問で、この範囲の基本的な区別を使用する。後の法改正を過去の登録要件へ遡及して当てはめない。

\subsection{自然法則の利用とソフトウェア}
\label{theory:compliance:invention}
自然法則自体の発見、抽象的な数学式、ゲームの約束、情報の単なる提示と、その法則を利用して問題を解く技術的な創作を区別する。特許と実用新案の定義はともに自然法則の利用を含む\parencite{jpo_patent_utility}。

ソフトウェアは記号を処理するから全て対象外、あるいはプログラムだから全て発明、とはいえない。請求項全体が自然法則を利用した技術的な創作であるかを確認する。抽象的な暗号・算法と、それを具体的な情報処理装置で実現する発明も区別する。特許庁審査基準はソフトウェアを含む全体の判断を示す\parencite{jpo_invention_criteria}。発明該当性と、新規性・進歩性等の充足は別の段階である。

\subsection{営業秘密の3要件}
\label{theory:compliance:secret}
秘密管理性、有用性、非公知性を満たす情報が営業秘密である\parencite{meti_trade_secret}。社内情報という名称だけで足りず、アクセス制御や表示等により秘密として管理する。特許のように公開と登録による独占を得る制度とも異なる。知的財産の創作物に含むかを問う問題と、産業財産権4種類に含むかを問う問題で、対象範囲を取り違えない。

\section{著作権の制限：許諾なしの利用にも要件がある}
\label{theory:compliance:citation}
引用は、自分の検討・説明に他者の表現を必要な範囲で取り込む利用である。公表された著作物について、公正な慣行に合致し、引用の目的上正当な範囲を満たすことを確認し、出所を示す。自分の記述と引用部分を区別し、利用目的・量・方法・影響を検討する。出典を付けた事実だけで長大な転載が適法になるわけではない\parencite{bunka_citation_faq}。

平成27年の設問でも引用を学術研究だけへ限定する説明は誤りである。適法な引用の範囲では翻訳利用も可能であるが、翻訳と自由な翻案・改変を同一視しない\parencite{kyushu_copyright2015}。公表は権利放棄を意味しない。

官公資料には、32条2項の別の仕組みもある。公的機関が周知のため自らの名義で公表する広報・調査統計・報告等を説明材料として刊行物へ転載する場合で、禁止表示などの条件を確認する\parencite{cabinet_copyright2008}。官公庁サイトにある全ての著作物が無条件に対象とはしない。

\subsection{非享受目的の利用とAI}
\label{theory:compliance:nonenjoyment}
通常の鑑賞・読書は表現された思想や感情を享受する。一方、複製機器の性能試験や機械的な情報解析は、表現の鑑賞とは異なる目的で利用する場合がある。2018年改正・2019年施行の30条の4は、この非享受目的について必要な限度で利用を認め、不当な利益侵害を除外する\parencite{bunka_amend2018}。

AI開発の学習データ収集や、その開発を行う第三者への提供が該当し得る。ただし既存作品の表現を出力させる目的が併存する利用や、解析用データベースの販売市場に衝突する利用などは要件を別に検討する。学習段階と生成物の利用段階の権利問題も別である\parencite{bunka_ai2024,bunka_text2026}。令和6年設問の例示を、AI利用一般の包括的な免許へ拡張しない。

\section{公正取引：独立した競争と市場情報の公正}
\label{theory:compliance:competition}
価格・数量を本来各事業者が独立に決める競争で、競争相手同士が共同決定すると、利用者は競争の便益を失う。価格・数量等のカルテル、入札の受注者・金額を事前に決める入札談合を、行為から識別する\parencite{jftc_rules}。入札談合は法的にはカルテルの典型例でもあり、用語の区別は概念が互いに完全に排他的という意味ではない\parencite{jftc_bid_cartel2021}。

\subsection{公共契約の例外と官製談合}
\label{theory:compliance:procurement}
国の契約は一般競争を原則とする一方、競争を許さない性質・目的、緊急等の法定条件では随意契約を認める。例外を適用する理由を具体的に説明し、事務の遅延だけを緊急性へ置き換えない。2006年の財務省通知もこの点を明確にする\parencite{mof_procurement2006}。

談合への行政措置、刑事責任、発注職員の関与に対する制度を区別する。官製談合防止法の職員の入札等妨害への刑事罰は2007年に施行され、平成27年出題時にも存在した\parencite{jftc_kansei}。各責任の構成要件は別に確認する。民間が発注するから競争法の対象外になることもなく、当時の公正取引委員会資料にも民間受注調整の処理実績がある\parencite{jftc_private2018}。

\subsection{金融市場の2種類の問題}
\label{theory:compliance:market}
インサイダー取引は、会社関係者等が法定の重要事実を職務等から知り、公表前に所定の取引をするという情報の公正の問題である。日常語の重要情報全般と法定の重要事実は同一ではなく、対象者・証券・公表・適用除外にも要件がある\parencite{fsa_insider_rules}。

相場操縦は、人為的に作った価格等を自然な需給の結果のように見せて他者の取引を誘引するなど、市場の価格形成を損なう問題である。大きな売買により価格が変わった事実だけで相場操縦と決まらない\parencite{fsa_market2004}。令和4・7年の問題では、競争事業者の共同決定、入札調整、非公開情報、偽った価格形成をそれぞれの名称へ対応させる。

\section{公益通報：問題を知らせる人を守る制度}
\label{theory:compliance:whistleblower}
事業者内部の違反を知る人が報復を恐れて沈黙すれば、事故や消費者被害は拡大し得る。公益通報制度は、対象事実を知らせる人の保護と法令遵守を結び付ける。単なる不満の表明や不正目的の告発、全ての機密の公開を無条件に保護する制度ではない。

\paragraph{出題当時：2018年}
公務員も原則として労働者に含む。内部通報、権限ある行政機関、その他外部の通報先によって要件が異なる。その他外部では相当の理由、被害防止に必要な通報先、報復や証拠隠滅のおそれ・内部調査が行われない・生命身体への急迫危険等の追加事情を確認する。追加事情は全ての同時成立ではなく法定の選択的条件である。改正前の消費者庁ハンドブックにより当時の範囲を確認する\parencite{caa_handbook_pre2022}。内部へ通報しない限り常に保護されないと一般化しない。

\paragraph{現在：2026年10月10日}
2020年改正は2022年6月1日から施行され、退職後1年以内の元労働者や役員等、対象・要件・体制を更新した。通報者の立場による要件の違いも確認する\parencite{caa_amend2020,caa_current_faq,caa_report_current}。300人以下の事業者への体制整備は努力義務であり、法全体が小規模事業者へ適用されないという意味ではない。

\paragraph{施行予定の改正}
2025年改正の主要部分の施行日は2026年12月1日である\parencite{caa_schedule2026}。この教材の確認日では将来の制度なので、施行中の制度や平成30年の正誤へ遡及して使わない。通報先と必要な情報を調べ、手続・確認状況を記録することは、倫理の情報開示の議論と接続する。

\paragraph{省略された条件を全リストにしない}
平成24年度の外部通報の4記述は、旧制度の報復、証拠隠滅、正当な理由のない通報禁止要求、
書面通報後20日以内の調査通知の欠如に対応する。
対象事実や相当の理由等も前提となり、4文だけで全ての要件を尽くすとはしない。\cite{saf53_whistle_oldfaq}
平成27年度の別表の略記は政令の指定を含めて読み、外部通報条件の代表列挙も排他的な全条件とはしない。
公益通報でも、他人の正当な利益や公共の利益を害さないよう配慮する。
第三者の個人情報を無制限に拡散する権利が生じるわけではない。\cite{caa_handbook_pre2022}

\paragraph{原典の改正年表現と対象者}
令和5・7年度の原典は2022年又は令和4年改正と記すが、2020年改正の2022年施行と区別する。
公務員も対象に含まれる。公務員法が適用されるから全て除外という肢を否定する。\cite{saf53_whistle_oldpublic,caa_amend2020}
300人以下の体制整備の努力義務と、通報者の法定保護は別である。
改正の成立、公布、施行、資料の確認日を一つの日付へまとめない。

\section{標準化と認証：規格を決めることと適合を示すこと}
\label{theory:compliance:standards}
規格は寸法、性能、試験方法等の共通の約束を与え、互換性と比較可能性を支える。ISOは産業分野、IECは電気・電子技術分野を主に扱う。JISCの基本説明ではISOの対象から電気・通信・電子技術を除くと整理するが、情報技術等には共同の標準化もある\parencite{jisc_isoiec}。

規格の制定と、具体的な製品等がそれに合うことの認証は別である。JISマークは国登録の第三者機関による審査・認証を受けて表示する。定期審査は少なくとも3年ごと、必要に応じ臨時審査や国の立入検査がある\parencite{jisc_certification}。マークを描くだけで適合性を証明することはできない。

\paragraph{変更の年を分ける}
新JISマーク制度は2005年10月に始まり、旧表示の経過措置は2008年9月までである\parencite{mhlw_jis_history}。2019年7月の改正では規格名が日本工業規格から日本産業規格、法名が工業標準化法から産業標準化法へ変わり、データ・サービス等へ対象が拡大した\parencite{meti_jis2019}。2019年をマークの新旧交替年としない。

令和6年設問の左右の説明は実図と逆である。現行の横並びJIS図はJISCの公式画像で確認できる\parencite{jisc_marks}。図の意味は文字の説明だけで決めず、原典そのものと独立資料を照合する。

\section{出題の読み方と本章の範囲}
正誤表は、各文の判断を先に書いてから行へ対応させる。数の問題は正しい文の数と選択肢番号を区別する。引用と複製、知的財産権と産業財産権、契約方式と競争制限、ISO等の規格と製品認証、改正の成立と施行を切り分ける。制度の専門的な適用には個別条件があり、ここでの試験要旨を例外を含む全規定の説明へ置き換えない。

以下はAstra査読前の独自解説初稿である。未執筆問題の一覧も残し、全問の確認を継続する。
\section{公衆と秘密保持：情報を伝える目的と範囲}
\label{theory:compliance:public_information}
憲法12条は権利の保持と濫用を避ける責任を、民法1条は私権と公共の福祉の適合を定める\parencite{jp_constitution,jp_civil}。技術者倫理では、技術の結果を受けても、その専門的内容を知って同意し又は回避する機会が乏しい人々への責任が重要になる。知識の非対称性は、個人の能力一般を低く評価する根拠ではない。

説明責任は情報を正確に伝え、根拠や不確実さへ応答する責任である。必ず相手が納得するという結果を強制できるわけではない。国家公務員法100条の職務上の秘密保持は退職後も続き、法に定める発表の許可等も区別する\parencite{jp_national_public}。説明責任と秘密保持をどちらか一語だけで決めず、情報の性質、法的根拠、開示先と目的を検討する。

\subsection{通報を是正へ結び付ける}
不正の根拠を検証し、相談・通報手段を活用し、適切な場合は上司へ異議を伝え、是正案を示す。上司の反応だけで終わらず他の理解者や経路を探すことも必要になる。専門職の支持を求めるのは事実を検討するためで、結論に合わせて口裏を合わせることではない。

根拠確認の努力を違法な証拠収集の許可へ広げない。内部の経路を活用するという倫理的助言を、外部通報の保護には必ず内部通報を先行しなければならないという法的条件へ変換しない。公益通報の保護は対象・目的・通報先等の要件で判断する。出題当時から、まさに生じようとしている対象事実を含む点は、後年の対象者拡大とは別である\parencite{caa_handbook_pre2022,caa_current_faq}。
\section{知的財産を分類するための二つの問い}
\label{theory:ip:scope}
知的財産の問題では、まず保護される対象を問い、次にどの制度で保護するかを問う。技術的な解決方法、文章や画像の表現、商品の出所を示す標識、秘密の製造情報は、同じ情報でも保護の根拠が異なる。情報は多人数が同時に利用できるため、物の占有と同じ方法で利用を制御できない。制度は創作・研究開発への誘因と成果の利用を調整する。

特許庁が産業財産権と総称するのは、特許権、実用新案権、意匠権、商標権の4種類である\parencite{ip_industrial_four}。著作権、回路配置利用権、育成者権は知的財産に関係する権利であっても、この4種類に含まれない。営業秘密、商号、商品等表示の保護も広い知的財産の体系に位置付く\parencite{ip_basic_act,ip_broad_scope}。広い体系へ含めることは、全てが登録により絶対的な独占権を生むという意味ではない。

\subsection{集合の補集合と選択肢番号}
\label{theory:ip:count}
列挙された対象の集合を \(L\)、4種類の集合を \(I\) とすると、産業財産権に含まれない対象は \(L\setminus I\) である。平成30年問6は広い知的財産の範囲外の数、令和2年問5は産業財産権の範囲外の数を問う。同じ名称が並んでも、引く集合が違えば数が変わる。対象数を求めた後で、その数を示す選択肢番号に対応させる。

\section{特許：定義と権利取得の要件を分ける}
\label{theory:ip:patent}
自然法則を利用した技術的思想の創作のうち高度なものが、特許法上の発明である。発明に該当しても、直ちに特許を得るわけではない。産業上利用できること、出願前の公知・公用・刊行物等に照らした新規性、進歩性などを別に審査する。公の秩序、善良の風俗又は公衆の衛生を害するおそれがある発明は保護対象から除かれる\parencite{ip_patent_statute}。

\subsection{出願前公開と登録後公開の区別}
\label{theory:ip:novelty}
制度が技術の公開と利用を促すことと、出願前に既に公表されていることを取得要件にすることは別である。通常、出願前公開は新規性に不利に働く。ただし公開された発明が例外なく永久に特許を得られない、とも断定しない。新規性喪失の例外には対象行為、期限、手続の条件がある。

平成30年改正では例外期間が6か月から1年へ延長され、2018年6月9日以後の出願への適用に公開日の経過措置も置かれた\parencite{ip_novelty_grace2018}。令和元年再試験の判断は、この当時すでに例外が存在したことと整合させる。例外は「刊行物で発表済みでなければならない」という要件を作らない。

\section{著作権：表現、創作性、権利の発生}
\label{theory:ip:copyright}
著作物は思想又は感情を創作的に表現したものである。学術的な事実や抽象的なアイデアそのものと、その説明文・図の具体的な表現を区別する。創作性は芸術的な傑作という評価を要求するものではないが、法の定義から削除されたわけでもない。

著作者人格権と著作権の享有に申請や審査を要求しない仕組みを無方式主義という\parencite{ip_copyright_statute}。特許庁へ登録しなければ著作権が発生しないという説明は、産業財産権の手続との混同である。権利の発生と、一定事項の登録制度が存在することも別である。

\subsection{人格的利益と財産的利用}
\label{theory:ip:personality}
公表権、氏名表示権、同一性保持権は著作者の人格的利益に関わる。複製、公衆送信などの利用を管理する財産的な権利と区別する。歌手が歌唱する立場にあることだけで、他者が作った歌詞を自由に変更する権限は生じない。利用許諾があることも、全ての改変への同意があることを当然には意味しない。

\subsection{個人の趣味と公開送信}
\label{theory:ip:public_web}
個人的な趣味であっても、公開Webページへの画像掲載は私的な複製と同じ範囲ではない。写真集の画像について、許諾や引用等の適法な根拠を欠く掲載を、趣味という目的だけで正当化しない。写真の権利と被写体に関わる別の権利も分けて考える。

\section{引用：出典表示だけでは条件を満たさない}
\label{theory:ip:quotation}
32条1項の引用は、公表された著作物を、公正な慣行と引用目的上正当な範囲に従って利用する仕組みである\parencite{ip_copyright_statute}。自分の論述に必要な部分を区別し、利用量と目的の関係を説明できることが重要になる。48条の出所明示も確認するが、出典を書けばどの量でも転載できるという制度ではない。

権利制限を満たす引用は、保護期間内であっても許諾なしにできる。保護期間満了後の利用、権利者の許諾に基づく利用、法が認める引用は異なる根拠である。「無断なら例外なく違法」という断定を避ける。

\subsection{官公資料の範囲と禁止表示}
\label{theory:ip:government}
法令など13条に掲げる対象が著作権の対象から除かれることと、官公庁が作成した全資料を自由に使えることは別である。32条2項は、国等が一般への周知を目的に自らの名義で公表した広報・調査統計・報告等を、説明材料として新聞紙、雑誌その他の刊行物へ転載する仕組みを定めるが、転載禁止の表示がある場合を除く。対象、目的、媒体、禁止表示を確認する\parencite{ip_copyright_statute}。

\subsection{引用のための翻訳と自由な翻案}
\label{theory:ip:translation}
適法な引用に対応する範囲では翻訳利用が認められ、外国語の原文をそのまま使うことだけが唯一の方法ではない。これは他者の文章を自由に翻案・改変してよいという意味ではない。取得した現在の条文では47条の6が翻訳等を扱う。平成25・26年度の問題について当時の条番号を現在番号に置き換えた証跡とはしない。二問の当時版全文の照合は未完として残す。

\section{営業秘密：情報の種類と法的成立を分ける}
\label{theory:ip:secret}
不正競争防止法2条6項の営業秘密は、秘密管理性、有用性、非公知性を同時に満たす情報である\parencite{ip_unfair_statute}。顧客名簿や新規事業計画はその典型になり得るが、資料名や社内で作成された事情だけで営業秘密になるわけではない。秘密として管理され、事業に有用で、公然と知られていないことが必要である。

\subsection{有用性と現実の使用}
\label{theory:ip:usefulness}
情報の有用性は、現に製品へ使われていることだけでは決まらない。将来の事業計画、失敗した実験の記録も、不要な研究費を避けるなどの価値を持ち得る。一方、脱税や有害物質の垂れ流し等の反社会的な活動の情報は、法が保護する正当な事業活動の有用性とは扱われない\parencite{ip_secret_handbook2017,ip_secret_commentary2018}。違法行為の情報を隠すことと、適法な技術情報を守ることを分ける。

\subsection{公開、秘密管理、漏洩経路}
\label{theory:ip:secrecy}
刊行物や公開された特許に記載された内容は、そのまま非公知の秘密とはできない。そこに記載されていない別のノウハウまで一緒に公開されたとは限らない。管理表示やアクセス制御は、対象者が秘密と認識できる管理の具体化である。漏洩経路には従業員、退職者、取引先、共同研究先、サイバー攻撃がある。経路の存在の確認と、特定年に急増したという量的主張の確認は別である。

\subsection{令和2年と令和3年の同じ省略表現}
\label{theory:ip:secret_ambiguity}
両年の設問は、顧客名簿等について、企業活動の過程で生まれたので営業秘密であると述べる。公式表は典型的な該当例として真とする組合せである。一方、三要件を省略しない一般命題として読むと、作成された由来だけでは成立を保証できない。令和3年の独立初回判断はこの読みから4、公式正答は5であった。

本教材は公式値を保持するとともに、初回判断と省略条件を残す。公式と一致するよう法の要件を変更しない。令和2年も同じ条件省略があるため、公式値と一致した事実だけで曖昧さを解消済みにしない。

\section{非享受目的：利用の目的から権利制限を判断する}
\label{theory:ip:nonenjoyment}
30条の4は、表現された思想又は感情を自ら享受し又は他人に享受させることを目的としない場合に、必要な限度で利用を認める。ただし著作権者の利益を不当に害する場合を除く。技術試験、情報解析、人の知覚によらない計算処理は例示であり、「AIだから許される」という機器名の判断ではない\parencite{ip_flexible2019,ip_copyright_statute}。

\subsection{六つの例とそれぞれの境界}
\label{theory:ip:six_cases}
2019年文化庁資料は、AI学習用の収集・第三者提供、プログラムの解析、複製機器や和紙の性能試験、送り仮名等の調査、場所の3DCG化、全文検索用の複製を具体例に挙げる\parencite{ip_flexible2019}。それぞれ表現の鑑賞ではなく、解析・性能・事実情報等を目的とする例である。

AI学習では、既存作品の表現を享受させる目的の併存や、解析用に販売されるデータベースの市場への影響を別に確認する。性能試験の複製物を鑑賞用に頒布する行為まで認めることにはならない。写真から場所を再現する場合も、元写真の創作的な表現を3DCGで承継して享受させる場合は区別する。

検索のための内部複製と検索結果の表示も別である。書名・著者名等の表示を目的とする例と、本文の一部を表示する利用についての47条の5の条件を混同しない。プログラム解析が著作権法上許される場合でも、契約など別の問題が全て解決するわけではない。

\section{法の時点を保った学習の手順}
\label{theory:ip:historical}
4種類への分類、引用の条件、三要件の論理は共通に参照し、保護対象の拡大や条番号、例外の改正は出題時点を確認する。非享受目的の2019年施行の規定を、それ以前の問題にそのまま遡及しない。現在の条文の説明、当時の行政資料で確認した範囲、未取得の過去版全文を区別する。

正誤表では各文の判断を先に定め、原典の列順で5候補と照合する。否定、数、組合せの違いを保ち、同じ公式番号が続くことを同じ答えの内容とみなさない。

\section{雇用と職場：専門能力を発揮できる条件を守る}
\label{theory:employment:background}
技術者は安全な製品を作る責任だけでなく、設計・製造・点検を担う人の尊厳と健康にも関わる。威圧によって異常を報告できない職場では、技術上の問題を発見しても対処へ結び付かない。長時間労働や特定個人への業務集中も判断・継承を損なう。雇用法令と職場の仕組みを、専門能力が継続して発揮されるための条件として理解する。

\subsection{労働法：主体と義務を先に読む}
\label{theory:employment:law}
伝統的な労働三法は労働基準法・労働組合法・労働関係調整法である\parencite{employment_labor_three}。最低条件の保護、集団的な交渉、紛争の調整という役割を区別する。労働安全衛生法は労働者の安全と健康を扱い、規模等に応じた産業医の選任を求める\parencite{employment_industrial_physician}。技術士資格があるだけで、労働者を使う際の義務が免除されるわけではない。

36協定は法定の時間外・休日労働に関する労使協定である。締結と届出を区別し、届出だけで上限を自由に撤廃できるものと読まない\parencite{employment_overtime}。以下では出題当時の記述を保ち、2026年10月12日時点で確認した変更は別に述べる。全ての告示の改正履歴を照合し終えたという意味ではない。

\subsection{パワハラ：職位より関係、叱責より態様}
\label{theory:employment:power}
平成24年WG報告は、職場内の優位性を背景に、適正な業務範囲を超えて苦痛又は環境悪化を生む行為を整理した。上司から部下だけでなく、知識や協力への依存による同僚・部下の優位性も含む。身体攻撃・精神攻撃・隔離・過大要求・過小要求・私生活への過度な介入は典型例であり、網羅的な六項目の暗記だけでは足りない\parencite{employment_power2012}。

検討では、目的に必要な指示か、方法が相当か、相手と周囲の環境へどう影響するかを分ける。手順の誤りを具体的に指摘することと、人前で人格を否定することは同じでない。主観的不満だけで該当性は決まらないが、相手の苦痛を調べずに「指導だから問題なし」としてもいけない。現在の指針も三要素と客観的な適正指導の区別を定める\parencite{employment_power_current}。

\subsection{セクハラ：不利益と環境悪化}
\label{theory:employment:sexual}
性的言動への対応で処遇上の不利益を受ける対価型と、言動によって仕事の環境が害される環境型を区別する。同性間でも生じ、労働者の雇用形態や行為者の所属だけで保護を除外できない。派遣労働者には派遣元・派遣先の関係を確認する\parencite{employment_sexual_current}。

拒絶を口にできない関係や、問題化すると仕事を失うと感じる状況があるため、沈黙だけから同意を推定しない。相談を受けた組織は事実を確認し、情報の扱いと相談者への不利益を防ぐ体制を整える。2018年研修資料も措置義務を説明しており、単なる努力目標という過去問の肢は誤る\parencite{employment_sexual2018}。

\subsection{各類型の範囲を同一にしない}
\label{theory:employment:scope}
日常語のハラスメントは広い。セクハラ、パワハラ、妊娠等に関するハラスメントは法的な条件が同じではない。顧客からの性的言動もセクハラの行為者範囲に入るが、その事実だけで顧客との全ての紛争を同じ職場内のパワハラ定義へ入れない。古いWG報告の顧客例は、職場内の定義とは別の脚注である。

性的指向や性自認を侮辱したり、本人の同意なく周囲へ伝えたりすることは、人格と機微情報に関わる。技術指導に必要な行為かを確認する。過去問の「本人はもとより職場全体のハラスメントとなり得る」は、周囲にも就業環境の悪化が及び得るという判断であり、全員が必ず同一の法的被害者となるという意味ではない。

\subsection{産休と妊婦への配慮}
\label{theory:employment:maternity}
産前産後休業は母体保護を目的とし、育児休業とは根拠と条件が異なる。産前は原則出産予定前6週間、多胎14週間で本人の請求が条件である。産後は原則8週間の就業禁止で、6週間経過後に本人が請求し医師が支障なしと認める業務には例外がある。正規・有期・派遣といった雇用形態で一律に除外しない\parencite{employment_maternity2014,employment_maternity_current}。

妊娠を理由に仕事から一方的に排除することと、健康や業務状況を考えて本人へ調整を相談することは異なる。配慮という名称だけで不利益を正当化せず、目的・本人の意向・法定措置を確認する。2018年の公的研修例も、残業の負担を相談して減らす配慮と、権利利用を妨げる強要を分けている。

\subsection{労働時間：場所と自己管理は免除理由でない}
\label{theory:employment:worktime}
使用者の指揮命令下の時間を考える。命令された業務上の学習は、学習という名称でも労働時間に含まれる。使用者は時間を適切に把握し、自己申告を用いる場合にも記録との食違いを確認する\parencite{employment_hours2017}。

管理監督者は肩書だけで決まらず、職務・権限・勤務の裁量・待遇等の実態から判断する。該当する場合も深夜割増賃金や年次有給休暇まで一律に除外されるものではない\parencite{employment_manager2008}。勤務終了後に必要な報告や研修を行わせながら記録から除外することは、健康管理と制度の趣旨を損なう。

\subsection{裁量労働：資格名と業務の列挙を分ける}
\label{theory:employment:discretion}
専門業務型裁量労働制は、法令で指定する業務の中から労使協定で対象を定め、定めた時間働いたものとみなす仕組みである。平成25年当時は19業務であり、保存した平成26年3月資料にも同じ19業務が並ぶ。この一覧に「技術士の業務」はない\parencite{employment_discretion2014}。

技術士が研究開発など対象業務に従事する可能性と、技術士資格を持てば自動的に対象になる命題は異なる。この混同で平成25年度問4の初回解答を誤ったため、初回3と資料確認後の訂正5を別々に保存した。適用を確かめるには肩書から推測せず、実際の従事業務と指定範囲を照合する。

\paragraph{現在の変更}2024年4月から対象は20業務で、本人同意・同意の撤回等の要件が加わった。現行一覧にも技術士業務という列挙はない\parencite{employment_discretion_current}。みなし時間は健康管理のための勤務状況把握を不要にするものではない。

\subsection{フレックスとテレワーク}
\label{theory:employment:flex}
フレックスは一定の清算期間の総労働時間を定め、その枠内で始業・終業を労働者が選択する仕組みである。期間の上限は2019年4月の改正で1か月から3か月へ延長された\parencite{employment_reform2018}。平成29年度の説明へ3か月を遡及しない。柔軟な時間配分と、総労働時間の把握を両立する。

\label{theory:employment:telework}
テレワークは場所の選択を広げるが、雇用型では労働法令が適用される。勤務場所を示し、時間を把握し、中抜けで自由利用が保障される時間と、移動中も命令された仕事をする時間を分ける。平成30年指針の概要を載せた令和2年白書でもこの区別を確認できる\parencite{employment_telework2020}。時間単位年休には別の制度要件があり、中抜けしたことだけで勝手に年休へ振り替えない。

現在の指針は記録・申告の方法、休止と移動、長時間化対策を整理する\parencite{employment_telework_current}。通信できる環境は常時応答の義務と同じでない。自己管理という説明によって、組織側の時間把握や負荷調整を放棄しない。

\subsection{働き方改革：上限・割増・休暇を分ける}
\label{theory:employment:reform}
法定の時間外労働上限は使用者の指示だけで緩和できない。年10日以上の年休が付与される労働者について、年5日の取得を確保する。本人の請求による取得や計画付与の取得分を含むため、全員へ別枠の5日を追加指定する説明は不正確である\parencite{employment_paidleave}。

月60時間超の時間外労働は割増賃金の問題であり、10時間ごとに代休1日という制度ではない。代替休暇には労使協定と本人意思を要し、引き上げた割増分に対応して換算する仕組みである\parencite{employment_alternative2008}。法定休日に働いた後の代休、事前に休日を振り替える操作、割増分の代替休暇は別の概念である。

\subsection{仕事と生活の調和：属人化と不要作業を減らす}
\label{theory:employment:worklife}
業務を標準化して引き継げる形にし、資料の過剰作成や探索を減らし、予定と負荷を共有する。これらは同じ仕事を少人数へ集中させるためでなく、時間と相互支援を確保する組織改善である。内閣府の十の実践は会議、資料、整理、標準化、時間、分担、相互理解、予定、集中、改善共有を扱う\parencite{employment_worklife2010}。

会議の到達点を明らかにする場合、技術上の証拠が不十分なら追加調査と責任者・期限を決めることも到達点になる。「必ず結論」という標語を、判断の根拠を省略して安全性を断定する命令へ変えない。効率化の取組後も、業務の質、負荷の偏り、実際の労働時間を観察し、改善の効果を確認する。

\section{安全保障貿易管理：仕様・相手・用途を確認する}
\label{theory:export:screening}
同じ測定器が民生研究にも軍事目的にも役立つ場合、買い手の肩書だけで用途は分からない。安全保障貿易管理は、軍事転用可能な貨物や技術が懸念活動へ渡ることを防ぐため、外為法等に基づいて行う\parencite{export_overview}。

まず貨物や技術の仕様を法令のリストへ照合する該非判定を行う。次に、相手・需要者・用途・仕向地等から取引を審査する。リスト非該当でも、補完的なキャッチオール規制の条件に該当する場合がある。最後に、審査した物と発送する物の同一性、許可の取得と内容を出荷時に確認する。少量サンプル、国内で買える製品、自作装置、手荷物での持出しというだけで、この確認を省けない。

メーカーの判定書は重要な資料だが、古い判定書が現在の型式・仕様・法令に合うかを輸出者も確かめる。信頼する取引先からの紹介も、新たな需要者の用途確認を代替しない。営業・技術・管理・出荷の各部門で判断根拠を共有することで、適切な審査を実際の発送へ結び付ける。

\section{技術提供と例外：場所より情報の内容を調べる}
\label{theory:export:exceptions}
技術には設計・製造・使用に必要な情報を、データや助言として提供する場合がある。国内の非居住者への提供も対象となり得るので、日本国内であることだけで規制外と判断しない\parencite{export_scope2022}。大学の非公開な共同研究打合せで技術を交換することと、不特定多数が入手できる公知技術の発表は違う。

装置の部分品特例も、部品が組み込まれているという一事だけで使えるわけではない。公的事例集は、ロボットに内蔵されたセンサーが主要要素でないことを学内の担当部署が確認する手順を示す\parencite{export_cases}。特例の対象と仕様・価格等の条件を検討し、本人だけの思い込みで適用しない。特許出願等に必要な最小限の技術提供についての例外も、手続目的を越えるノウハウ提供まで含めない。当時条文の直接照合は初稿の留保事項である。

\begin{tcolorbox}[title=出題時点と現在の規則]
ここでの過去問は2021--2023年の事例である。経済産業省の現在の案内には2025年10月9日施行の補完的規制見直しが含まれる\parencite{export_current_updates}。現在の品目・対象国・用途・通知条件を過去問へ補って判定しない。取引の実施では、実施時点の規則を確認する。
\end{tcolorbox}

\section{専門職・継続管理・要求性能を別の軸に置く}\label{theory:compliance:management}
高度な専門能力を持つ人材の養成を目的とする専門職大学院の課程は、
博士後期課程そのものではない。専門職の学位と研究の博士課程を区別する。
法務博士という学位名の「博士」だけから博士後期課程と同一視しない。\cite{eng52_professional}

BCMは障害に備える計画を作り、訓練・点検・見直しを続けて事業継続能力を管理する活動、
BCPはその事業継続計画である。計画書を置くだけで管理の全過程が完了するわけではない。
当時の理解は平成21年版等の資料から確認し、後年の制度改訂を遡って加えない。\cite{eng52_bcp}

性能の要求は、果たすべき機能・安全・使用性等を定める。
製作仕様は材料・寸法・方法等を指定する手段であり、両者を同じ定義にしない。
性能規定で要求の満足を検証することと、仕様を守ればよいという評価を区別する。
現実の要求水準には性能と仕様が併存する場合もあるので、一方だけが常に許されるとはしない。
信頼性は必要な機能を条件・時間の中で遂行する度合いを確率的に扱う別の軸である。\cite{eng52_performance}

\subsection{VEの価値と品質の最低条件}
VEは必要な機能を確実に果たし、ライフサイクルを含む資源・コストと機能を比較して価値を改善する。
説明用の比\(V=F/C\)では、分母を下げるだけで機能の必要水準を失う変更を価値向上とはしない。
安い価格と、設計・保守・更新を含む全体コストを分ける。\cite{eng52_ve}

\subsection{内部統制の四目的と時点}
平成19年の基準では、業務の有効性・効率性、財務報告の信頼性、
事業に関わる法令等の遵守、資産の保全という四目的が示される。
組織の全員が業務の中で遂行するプロセスで、互いに関連する目的である。
経営者の整備責任と、その他の者が運用に関与することは両立する。
内部統制全体の目的と、金融商品取引法の財務報告に係る評価・報告・監査の対象を分ける。\cite{eng52_control}
平成23年度の表現を、後年の報告範囲の変更後の用語に無断置換しない。

\section{食品表示：どの情報が選択と安全を支えるか}
\label{theory:compliance:foodlabel}
期限、原産国・原料原産地、成分・アレルギー、事業者情報は、それぞれ別の判断を支える。
賞味期限は定めた保存条件のもとで期待される品質が保たれる期間、
消費期限は安全性を欠くおそれがないと認められる期間である。
いずれも保存条件が前提で、開封後や不適切な温度に置いた食品の無条件保証ではない。
賞味期限を過ぎたら直ちに全て食べられない、という安全期限との混同を避ける。\cite{saf53_food_expiry,saf53_food_loss2013}

輸入加工食品の最終加工国と、国内で加工した食品に使う原料の産地を分ける。
平成25年度の品目・原料の条件を、後年の広い対象へ変えない。
食品表示法の公布2013年と施行2015年も区別する。
アレルギー表示の指定は見直されるため、当時の指定品目を現在の数に書き換えない。\cite{saf53_food_notice}
表示は選択だけでなく、事故時の対象製品と責任を追跡する情報でもある。

\section{不実証広告：表示を裏付ける二つの条件}
\label{theory:compliance:adevidence}
効果を示す資料があるという事実だけでは、広告の主張は確かめられない。
資料自体の客観性と、表示された対象・条件・効果への対応の両方が必要である。
別の材料の試験や特殊条件の結果を、製品の一般的な利用へそのまま拡張しない。
行政が資料を求めても自由に断れるという制度ではなく、所定の不提出には制度上の効果がある。\cite{saf53_ad_overview}

自社試験でも妥当な方法で実証できるが、実施者が専門家であるだけでは足りない。
専門分野で一般に認められない特異な見解や、新しい技術を称することによって
試験・調査の実証を免れるものではない。
経験則を根拠とする場合も、存在を裏付ける専門的な見解等を確認する。\cite{saf53_ad2016}
2016年の説明では資料提出は原則15日で、新しい試験が必要という理由だけで延長できるとはしない。
2014年の設問へ現行条番号や後年の処分制度を補わず、当時の枠組みを保存する。

\section{PDCA：点検した結果を改善へ戻す}
\label{theory:management:pdca}
計画し、実施し、結果を点検し、是正と改善へ反映する順序がPDCAである。目標と評価方法を定め、実行の記録を取り、目標との差と原因を調べて次の計画へつなぐ\parencite{eth54_pdca2017}。

Checkは単に点検作業をしたという事実ではなく、計画に照らして結果と過程を評価する段階である。Actも行動一般ではなく、その結果を踏まえる是正・改善を表す。研究開発の評価は次の研究や資源配分へ役立てる。事前・中間・事後評価の時点と、PDCAの語順は別々に読む。

\subsection{平成23年度 問13}
\label{q:h23-reg-apt-13}

\noindent\textbf{原典：} sources/H23/aptitude.pdf、PDF物理14頁。

\textbf{公式正答：} 3（single）。

\par\textit{原典本文の校正と独自解説は未作成。分類のみの登録項目。}

\subsection{平成23年度 問15}
\label{q:h23-reg-apt-15}

\noindent\textbf{原典：} sources/H23/aptitude.pdf、PDF物理16頁。

\textbf{公式正答：} 1（single）。

\par\textit{原典本文の校正と独自解説は未作成。分類のみの登録項目。}

\subsection{平成27年度 問10}
\label{q:h27-reg-apt-10}

\noindent\textbf{原典：} sources/H27/aptitude.pdf、PDF物理10頁。

\textbf{公式正答：} 5（single）。

\par\textit{原典本文の校正と独自解説は未作成。分類のみの登録項目。}

\subsection{平成27年度 問12}
\label{q:h27-reg-apt-12}

\noindent\textbf{原典：} sources/H27/aptitude.pdf、PDF物理12頁。

\textbf{公式正答：} 2（single）。

\par\textit{原典本文の校正と独自解説は未作成。分類のみの登録項目。}

\subsection{平成28年度 問6}
\label{q:h28-reg-apt-06}

\noindent\textbf{原典：} sources/H28/aptitude.pdf、PDF物理6頁。

\textbf{公式正答：} 1（single）。

\par\textit{原典本文の校正と独自解説は未作成。分類のみの登録項目。}

\subsection{平成28年度 問11}
\label{q:h28-reg-apt-11}

\noindent\textbf{原典：} sources/H28/aptitude.pdf、PDF物理11頁。

\textbf{公式正答：} 4（single）。

\par\textit{原典本文の校正と独自解説は未作成。分類のみの登録項目。}

\subsection{平成29年度 問13}
\label{q:h29-reg-apt-13}

\noindent\textbf{原典：} sources/H29/aptitude.pdf、PDF物理12頁。

\textbf{公式正答：} 3（single）。

\par\textit{原典本文の校正と独自解説は未作成。分類のみの登録項目。}

\subsection{平成30年度 問5}
\label{q:h30-reg-apt-05}

\noindent\textbf{原典：} sources/H30/aptitude.pdf、PDF物理5頁。

\textbf{公式正答：} 4（single）。

\par\textit{原典本文の校正と独自解説は未作成。分類のみの登録項目。}

\subsection{平成30年度 問8}
\label{q:h30-reg-apt-08}

\noindent\textbf{原典：} sources/H30/aptitude.pdf、PDF物理8頁。

\textbf{公式正答：} 3（single）。

\par\textit{原典本文の校正と独自解説は未作成。分類のみの登録項目。}

\subsection{令和4年度 問9}
\label{q:r04-reg-apt-09}

\noindent\textbf{原典：} sources/R04/aptitude.pdf、PDF物理8頁。

\textbf{公式正答：} 5（single）。

\par\textit{原典本文の校正と独自解説は未作成。分類のみの登録項目。}

\subsection{令和4年度 問12}
\label{q:r04-reg-apt-12}

\noindent\textbf{原典：} sources/R04/aptitude.pdf、PDF物理11頁。

\textbf{公式正答：} 5（single）。

\par\textit{原典本文の校正と独自解説は未作成。分類のみの登録項目。}

\subsection{令和5年度 問4}
\label{q:r05-reg-apt-04}

\noindent\textbf{原典：} sources/R05/aptitude.pdf、PDF物理4頁。

\textbf{公式正答：} 2（single）。

\par\textit{原典本文の校正と独自解説は未作成。分類のみの登録項目。}

\subsection{令和6年度 問3}
\label{q:r06-reg-apt-03}

\noindent\textbf{原典：} sources/R06/aptitude.pdf、PDF物理3頁。

\textbf{公式正答：} 1（single）。

\par\textit{原典本文の校正と独自解説は未作成。分類のみの登録項目。}

\subsection{令和6年度 問12}
\label{q:r06-reg-apt-12}

\noindent\textbf{原典：} sources/R06/aptitude.pdf、PDF物理12頁。

\textbf{公式正答：} 3（single）。

\par\textit{原典本文の校正と独自解説は未作成。分類のみの登録項目。}

\subsection{令和7年度 問11}
\label{q:r07-reg-apt-11}

\noindent\textbf{原典：} sources/R07/aptitude.pdf、PDF物理11頁。

\textbf{公式正答：} 4（single）。

\par\textit{原典本文の校正と独自解説は未作成。分類のみの登録項目。}

